% Sparse Multi-Resolution 3D Diffusion: Comprehensive Survey
% For RDD-compatible 3D Generation with Voxel-like Representations
% Compile: pdflatex -interaction=nonstopmode sparse_3d_diffusion_survey.tex (run twice for TOC)
\documentclass[11pt, a4paper]{article}
\usepackage[utf8]{inputenc}
\usepackage[T1]{fontenc}
\usepackage{amsmath, amssymb, amsthm, mathtools}
\usepackage{geometry}
\geometry{margin=2.0cm}
\usepackage{booktabs}
\usepackage{longtable}
\usepackage{multirow}
\usepackage{hyperref}
\hypersetup{colorlinks=true, linkcolor=blue, citecolor=blue, urlcolor=blue}
\usepackage{listings}
\usepackage{xcolor}
\usepackage{enumitem}
\usepackage{float}
\usepackage{caption}
\usepackage{array}
\usepackage{tabularx}

\lstset{
    basicstyle=\ttfamily\small,
    keywordstyle=\color{blue}\bfseries,
    commentstyle=\color{green!60!black},
    stringstyle=\color{red!70!black},
    numbers=left, numberstyle=\tiny\color{gray},
    frame=single, breaklines=true, tabsize=4,
    showstringspaces=false, language=Python, captionpos=b
}
\lstdefinestyle{flowchart}{
    basicstyle=\ttfamily\footnotesize, frame=single, numbers=none,
    breaklines=false, language={}, keywordstyle={},
    commentstyle={}, stringstyle={}, backgroundcolor=\color{gray!5}
}

\newcommand{\alphabar}{\bar{\alpha}}

\title{\huge Sparse Multi-Resolution 3D Diffusion\\[0.4em]
\Large A Comprehensive Survey of 3D Representations\\Compatible with RDD's Constant Upsampling\\[0.2em]
\normalsize NeurIPS/ICLR Candidate Research Direction}
\author{Research Notes -- YANG, YUAN}
\date{July 2026}

\begin{document}
\maketitle
\tableofcontents
\newpage

%=================================================================
\section{Executive Summary}
%=================================================================

\subsection{Core Thesis}

RDD (Residual Drift Diffusion) introduces a \textbf{constant upsampling} mechanism in multi-resolution diffusion that does not depend on neighboring values. This property enables a \textbf{sparse diffusion} paradigm: voxels identified as empty at coarse resolution can be \emph{discarded} during refinement, since their upsampled values are independent of neighbors. This survey identifies 3D representations that can be mapped to $N \times N \times N$ grids and are compatible with this sparse multi-resolution diffusion framework.

\subsection{Key Finding: Top 3 Representations for RDD-3D}

\begin{enumerate}
    \item \textbf{GaussianCube} (NeurIPS 2024): 3DGS $\to$ Optimal Transport $\to$ $32^3 \times 14$ grid. Perfect grid structure, proven diffusion compatibility, extremely compact.
    \item \textbf{XCube/VDB Sparse Voxel Hierarchy} (CVPR 2024): Native sparse voxel octree with hierarchical latent diffusion. Scales to $1024^3$. Perfect for sparse RDD.
    \item \textbf{K-Planes / TensoRF-VM factorization}: Compact plane-based representations where upsampling operates on 2D factors. Implicit sparsity via Hadamard product.
\end{enumerate}

\subsection{Project Architecture Vision}

\begin{lstlisting}[style=flowchart, caption={Proposed Sparse RDD-3D Pipeline}]
                    RDD Sparse 3D Generation Pipeline
    ================================================================
    
    TRAINING:
    3D Asset --> Encoder (VAE) --> Latent Grid N^3 x C
                                        |
                                   [Sparsify: mark empty voxels]
                                        |
                                   Sparse Latent (only surface voxels)
                                        |
                            RDD Forward Process (multi-scale drift)
                                        |
                            Train Denoiser (Sparse 3D U-Net / DiT)
    
    GENERATION (Multi-Resolution):
    Step 1: Generate coarse grid (8^3 or 16^3) -- DENSE
            |
    Step 2: RDD Constant Upsample to 32^3
            |-- Identify empty voxels (density < threshold)
            |-- DISCARD empty voxels (sparse!)
            |-- Denoise only non-empty voxels
            |
    Step 3: RDD Constant Upsample to 64^3
            |-- Further sparsify
            |-- Denoise sparse subset
            |
    Step 4: RDD Constant Upsample to 128^3 (final)
            |-- Finest detail generation
            |
    DECODE: Sparse Latent --> 3D Output (mesh/3DGS/NeRF)
\end{lstlisting}

%=================================================================
\section{Background: RDD's Constant Upsampling Property}
%=================================================================

\subsection{Why Constant Upsampling Enables Sparse 3D}

In standard image super-resolution or cascaded diffusion, upsampling (bilinear, bicubic) \emph{depends on neighboring pixels}. This means every pixel must exist for correct interpolation. RDD's key innovation:

\begin{equation}
    \boxed{\text{UP}_{\text{RDD}}(x_i) = x_i \quad \text{(constant/nearest-neighbor replication)}}
\end{equation}

Each voxel's upsampled value depends \textbf{only on itself}, not neighbors. Consequences:
\begin{itemize}
    \item If a voxel is identified as ``empty'' (e.g., $\sigma < \epsilon$), we can skip upsampling it entirely
    \item The remaining non-empty voxels can be upsampled independently
    \item This converts the $O(N^3)$ dense computation to $O(k)$ where $k \ll N^3$ is the number of occupied voxels
    \item For typical 3D shapes, surface voxels scale as $O(N^2)$, giving $O(N^2/N^3) = O(1/N)$ sparsity ratio
\end{itemize}

\subsection{RDD Forward Process (3D Extension)}

\begin{equation}
    q(x_t | x_0) = \mathcal{N}\!\left(\sqrt{\alphabar_t}\, x_0 + \sqrt{\alphabar_t}\, \kappa(t)\, \mathbf{d},\; (1-\alphabar_t)\,\mathbf{I}\right)
\end{equation}

For 3D voxel grids:
\begin{equation}
    \mathbf{d}[i,j,k] = \mathrm{UP}(\mathrm{DOWN}(x_0))[i,j,k] - x_0[i,j,k]
\end{equation}

where $\mathrm{DOWN}$: $N^3 \to (N/r)^3$ (average pooling) and $\mathrm{UP}$: $(N/r)^3 \to N^3$ (nearest-neighbor). Since UP is constant:
\begin{equation}
    \mathbf{d}[i,j,k] = \bar{x}_0[\lfloor i/r \rfloor, \lfloor j/r \rfloor, \lfloor k/r \rfloor] - x_0[i,j,k]
\end{equation}

\textbf{Key insight}: $\mathbf{d}[i,j,k]$ only depends on $x_0[i,j,k]$ and its block average. If the entire block is zero (empty voxels), then $\mathbf{d} = 0$ and the forward/reverse process trivially produces zeros --- justifying discarding empty regions.

%=================================================================
\section{Category 1: Sparse Voxel Diffusion Methods}
%=================================================================

\subsection{XCube (CVPR 2024 Highlight) -- NVIDIA}

\textbf{Paper}: ``XCube: Large-Scale 3D Generative Modeling using Sparse Voxel Hierarchies''\\
\textbf{ArXiv}: \url{https://arxiv.org/abs/2312.03806}

\subsubsection{3D Representation: Sparse Voxel Hierarchy (SVH)}

Based on \textbf{VDB} (Volumetric Dynamic B+tree) --- the industry standard for sparse volumetric data:
\begin{itemize}
    \item 4-level shallow-wide tree: Hash Table $\to$ InternalNode1 ($32^3$ branching) $\to$ InternalNode2 ($16^3$) $\to$ LeafNode ($8^3$ voxels)
    \item $O(1)$ average random access per voxel
    \item Only allocates memory for non-empty regions
    \item Hierarchy: $\mathcal{G} = \{G_1, G_2, ..., G_L\}$ from coarse to fine
    \item Effective resolution up to $1024^3$ with only millions of active voxels
\end{itemize}

\subsubsection{Architecture}

\textbf{Stage 1: Sparse Structure VAE}
\begin{itemize}
    \item \textbf{Encoder}: Sparse 3D convolution network (fVDB/NanoVDB backend), downsamples $(G, A)$ to latent resolution $X$
    \item \textbf{Decoder}: Neural Kernel-inspired structure prediction. Predicts subdivision masks to progressively refine from $X$ back to target resolution
    \item Latent dimension: 8
    \item KL weight: 0.03
\end{itemize}

\textbf{Stage 2: Hierarchical Voxel Latent Diffusion}

\begin{lstlisting}[style=flowchart, caption={XCube Hierarchical Generation}]
Level 1 (16^3 DENSE):
    Diffusion Model: Dense 3D U-Net
    Output: Coarse structure (which regions are occupied)
        |
Level 2 (128^3 SPARSE):
    Condition: Level 1 features (via concat/cross-attn)
    Diffusion Model: Sparse 3D U-Net
    Output: Medium-detail geometry + attributes
        |
Level 3 (512^3 or 1024^3 SPARSE):
    Condition: Level 2 features
    Diffusion Model: Sparse 3D U-Net
    Output: Fine-detail surface + attributes
\end{lstlisting}

\subsubsection{Key Mathematical Details}

VAE loss:
\begin{equation}
    \mathcal{L}_{\text{VAE}} = \mathbb{E}_{q}[\log p(G,A|X)] - \beta \cdot D_{\text{KL}}(q(X|G,A) \| p(X))
\end{equation}

Diffusion loss per level (independent models):
\begin{equation}
    \mathcal{L}_{\text{diff}}^{(l)} = \mathbb{E}_{t,\epsilon}\left[\|\epsilon - \epsilon_\theta(x_t^{(l)}, t, c^{(l-1)})\|^2\right]
\end{equation}

where $c^{(l-1)}$ is the coarser level's output serving as conditioning.

\subsubsection{Experimental Results}

\begin{itemize}
    \item \textbf{ShapeNet}: 1-NNA (CD/EMD) competitive or superior to all baselines on Chair/Airplane/Car
    \item \textbf{Waymo Open Dataset}: 100m$\times$100m outdoor scenes, 10cm voxel size ($1024^3$ effective)
    \item \textbf{Generation speed}: $<30$ seconds for 1024$^3$ with millions of voxels
    \item \textbf{Framework efficiency}: 3$\times$ faster than TorchSparse, 0.5$\times$ memory
    \item Applications: unconditional, text-to-3D, scan completion, user-guided editing
\end{itemize}

\subsubsection{RDD Compatibility Analysis}

\begin{center}
\begin{tabular}{ll}
\toprule
Aspect & Compatibility \\
\midrule
Grid Structure & $\star\star\star\star\star$ --- native $N^3$ sparse grid \\
Constant Upsampling & $\star\star\star\star\star$ --- VDB supports arbitrary upsampling \\
Sparsity & $\star\star\star\star\star$ --- core design principle \\
Multi-resolution & $\star\star\star\star\star$ --- hierarchical by design \\
\midrule
\textbf{Limitation} & Uses cascade (independent models per level) \\
\textbf{RDD Advantage} & RDD's single-chain drift avoids cascade complexity \\
\bottomrule
\end{tabular}
\end{center}

%-------------------------------------------------------------------
\subsection{OctFusion (2024)}

\textbf{Paper}: ``OctFusion: Octree-based Diffusion Models for 3D Shape Generation''\\
\textbf{ArXiv}: \url{https://arxiv.org/abs/2408.14732}

\subsubsection{Core Innovation: Unified Multi-Scale U-Net}

Unlike XCube's cascade approach, OctFusion uses a \textbf{single unified U-Net} that simultaneously processes all octree levels:

\begin{itemize}
    \item No independent models per resolution
    \item Weight sharing across octree depths (depth 6=64$^3$, 7=128$^3$, 8=256$^3$)
    \item Feature propagation through octree neighborhood relations
    \item Only \textbf{33M parameters} (extremely lightweight)
    \item Generation: 50 steps, 2.5 seconds on single RTX 4090
\end{itemize}

\subsubsection{Mathematical Formulation}

Multi-scale diffusion with shared time step:
\begin{equation}
    x_t^{(l)} = \sqrt{\alphabar_t} \cdot x_0^{(l)} + \sqrt{1-\alphabar_t} \cdot \epsilon^{(l)}, \quad \forall l \in \{1,...,L\}
\end{equation}

Unified loss:
\begin{equation}
    \mathcal{L} = \sum_{l=1}^{L} w_l \cdot \mathbb{E}_{t,\epsilon}\left[\|\epsilon^{(l)} - \epsilon_\theta(x_t^{(l)}, t)\|^2\right]
\end{equation}

VAE (Dual Octree Graph Convolution):
\begin{equation}
    \mathcal{L}_{\text{VAE}} = \mathcal{L}_{\text{BCE}}(\text{occupancy}) + \mathcal{L}_{\text{MSE}}(\text{SDF}) + \lambda_{\text{KL}} \cdot \mathcal{L}_{\text{KL}}
\end{equation}

\subsubsection{Results}

\begin{itemize}
    \item Best average shading-image-based FID across all ShapeNet categories
    \item Particularly strong on complex structures (rifle, airplane)
    \item Supports: unconditional, category-conditional, text-conditional, sketch-conditional
    \item Output is guaranteed \textbf{continuous and manifold}
\end{itemize}

\subsubsection{RDD Compatibility}

\textbf{Extremely compatible} --- OctFusion's unified multi-scale approach is the \emph{exact same philosophy} as RDD:
\begin{itemize}
    \item Single model handling all resolutions (vs XCube's cascade)
    \item Shared weights = shared computation across scales
    \item RDD's $\kappa(t)$ schedule can naturally control which octree levels are active
    \item Direct architectural reference for RDD-3D
\end{itemize}

%-------------------------------------------------------------------
\subsection{TRELLIS (2024) -- Microsoft Research}

\textbf{Paper}: ``Structured 3D Latents for Scalable and Versatile 3D Generation''\\
\textbf{ArXiv}: \url{https://arxiv.org/abs/2412.01506}

\subsubsection{SLAT: Structured LATents}

\begin{itemize}
    \item Local latent vectors defined on \textbf{sparse voxel structure} (only active/surface cells)
    \item Sparse VAE encoder: 3D shifted-window attention transformer
    \item Serialization of occupied cells into active-voxel tokens
    \item KL regularization toward standard normal
\end{itemize}

\subsubsection{Two-Stage Pipeline}

\begin{enumerate}
    \item Generate sparse structure (which voxels are active) --- binary occupancy
    \item Generate latent vectors for non-empty cells --- continuous features
\end{enumerate}

\subsubsection{Architecture: Rectified Flow Transformer}

\begin{itemize}
    \item Backbone: Rectified Flow (straight-line ODE paths)
    \item Sparse attention via 3D shifted-window
    \item Model scale: trained up to \textbf{2 billion parameters}
    \item Training data: large-scale curated 3D assets
\end{itemize}

\subsubsection{Results}

\begin{itemize}
    \item TRELLIS.2: 512$^3$ in 3s, 1024$^3$ in 17s, 1536$^3$ in 60s (H100)
    \item State-of-the-art quality in both geometry and texture
    \item Supports text-conditional and image-conditional generation
\end{itemize}

\subsubsection{RDD Compatibility}

\textbf{Highly compatible}:
\begin{itemize}
    \item SLAT's ``sparse local latents'' directly maps to RDD's sparse grid
    \item Two-stage can be unified into RDD's multi-resolution framework
    \item Rectified Flow + sparse attention is a modern backbone option for RDD
    \item The shifted-window attention provides local-to-global information flow
\end{itemize}

%-------------------------------------------------------------------
\subsection{Direct3D-S2 (2025)}

\textbf{Paper}: ``Gigascale 3D Generation Made Easy with Spatial Sparse Attention''\\
\textbf{ArXiv}: \url{https://arxiv.org/abs/2505.17412}

\subsubsection{Spatial Sparse Attention (SSA)}

Key contribution --- efficient attention for sparse 3D:
\begin{itemize}
    \item \textbf{3.9$\times$ forward speedup} over dense attention
    \item Only 8 GPUs needed for 1024-resolution training
    \item 24-layer DiT blocks, hidden dimension 1024
    \item Grouped-Query Attention (GQA), group=2, 16 heads per group
    \item Operates only on non-empty voxel tokens
\end{itemize}

\subsubsection{RDD Compatibility}

SSA can be directly adopted as RDD's backbone attention mechanism for the sparse denoising network.

%=================================================================
\section{Category 2: Grid-Structured 3D Representations}
%=================================================================

\subsection{GaussianCube (NeurIPS 2024) -- Critical Method}

\textbf{Paper}: ``GaussianCube: Structuring Gaussian Splatting using Optimal Transport for 3D Generative Modeling''\\
\textbf{ArXiv}: \url{https://arxiv.org/abs/2403.19655}

\subsubsection{Method: 3DGS $\to$ Structured $N^3$ Grid}

\textbf{Step 1: Density-Constrained Fitting}

Standard 3DGS optimization but with a hard cap on Gaussian count:
\begin{equation}
    N_{\text{Gaussians}} \leq N_{\max} = N^3 \quad (N=32 \Rightarrow 32768 \text{ Gaussians})
\end{equation}

After each densification step, if count exceeds $N_{\max}$, prune lowest-opacity Gaussians. Train for 30K iterations.

Each Gaussian $G_i$ has 14 attribute channels:
\begin{equation}
    G_i = (\underbrace{\mu_i}_{\mathbb{R}^3}, \underbrace{s_i}_{\mathbb{R}^3}, \underbrace{q_i}_{\mathbb{R}^4}, \underbrace{\alpha_i}_{\mathbb{R}^1}, \underbrace{c_i}_{\mathbb{R}^3}) \in \mathbb{R}^{14}
\end{equation}

\textbf{Step 2: Optimal Transport Structuralization}

Solve the Linear Assignment Problem (LAP):
\begin{equation}
    \pi^* = \arg\min_{\pi \in \Pi_{N^3}} \sum_{i=1}^{N^3} \|\mu_i - g_{\pi(i)}\|_2^2
\end{equation}

where $g_j$ are predefined voxel center coordinates on the $N \times N \times N$ grid. Uses Jonker-Volgenant algorithm with segmented approximation (sort by coordinate, split into 4 segments, solve each independently).

\textbf{Step 3: Final Representation}
\begin{equation}
    \boxed{\text{GaussianCube} \in \mathbb{R}^{N \times N \times N \times 14}}
\end{equation}

Each voxel stores the full attributes of the Gaussian assigned to that position. Position channel stores \textbf{offset} $\Delta\mu = \mu - g_{\pi(i)}$ (residual from grid center).

\subsubsection{Diffusion Model}

Standard ADM 3D U-Net:
\begin{itemize}
    \item Replace all 2D ops (conv, attention, up/downsample) with 3D versions
    \item Full attention at $8^3$ and $4^3$ resolutions
    \item Input/output: $32 \times 32 \times 32 \times 14$
    \item 1000-step diffusion, cosine noise schedule
    \item Trained on 16 Tesla V100
\end{itemize}

\subsubsection{Experimental Results}

\begin{table}[H]
\centering
\caption{GaussianCube vs. Alternative Representations (ShapeNet Car)}
\begin{tabular}{lccccc}
\toprule
Method & PSNR$\uparrow$ & LPIPS$\downarrow$ & SSIM$\uparrow$ & Speed & Params \\
\midrule
Voxel $128^3\times32$ & 26.77 & 0.1103 & 0.9431 & 1$\times$ & 67.1M \\
Triplane $256^2\times32$ & 29.41 & 0.0721 & 0.9620 & 1.27$\times$ & 6.29M \\
\textbf{GaussianCube} $32^3\times14$ & \textbf{34.94} & \textbf{0.0347} & \textbf{0.9863} & 3.33$\times$ & \textbf{0.46M} \\
\bottomrule
\end{tabular}
\end{table}

Key findings:
\begin{itemize}
    \item \textbf{146$\times$ fewer parameters} than dense voxel, yet much higher quality
    \item \textbf{14$\times$ fewer parameters} than triplane
    \item Drastically better PSNR (+8.17 dB over voxel, +5.53 dB over triplane)
    \item FID scores superior on ShapeNet and OmniObject3D
\end{itemize}

\subsubsection{RDD Compatibility Analysis}

\begin{center}
\fbox{\parbox{0.9\textwidth}{
\textbf{GaussianCube is perhaps the single most compatible representation for RDD-3D:}
\begin{enumerate}
    \item Perfect $N^3$ grid structure --- standard 3D convolutions apply directly
    \item Each voxel has opacity $\alpha$ --- natural sparsity signal (discard $\alpha < \epsilon$)
    \item Constant upsampling: $16^3 \to 32^3 \to 64^3$ by replicating each voxel to $2\times2\times2$ children
    \item No neural decoder needed --- directly renderable via splatting
    \item Proven to work with standard 3D U-Net diffusion
    \item Extremely compact (460K params/shape)
\end{enumerate}
}}
\end{center}

\textbf{Proposed RDD-GaussianCube pipeline}:
\begin{lstlisting}[style=flowchart]
RDD Multi-Resolution GaussianCube Generation:

    t=T (noise): Start from 8^3 x 14 dense noise
        |
    RDD reverse (coarse stage): Denoise to clean 8^3 x 14
        |  (kappa transitions from 1->0 in this stage)
        |
    Constant Upsample: 8^3 -> 16^3 (replicate each voxel 2x2x2)
        |  Identify empty: alpha_i < threshold
        |  Discard empty voxels (SPARSE!)
        |
    RDD reverse (medium stage): Denoise 16^3 sparse
        |
    Constant Upsample: 16^3 -> 32^3
        |  Further sparsify
        |
    RDD reverse (fine stage): Denoise 32^3 sparse
        |
    Output: Sparse GaussianCube -> Render via splatting
\end{lstlisting}

%-------------------------------------------------------------------
\subsection{Plenoxels / Direct Voxel Representation}

\textbf{Paper}: ``Plenoxels: Radiance Fields without Neural Networks'' (CVPR 2022)\\
\textbf{ArXiv}: \url{https://arxiv.org/abs/2112.05131}

\subsubsection{Representation}

Each voxel corner stores:
\begin{equation}
    v_{ijk} = (\sigma_{ijk}, \mathbf{k}_{ijk}) \in \mathbb{R}^{1+27}
\end{equation}

where $\sigma$ is density and $\mathbf{k} \in \mathbb{R}^{27}$ are degree-2 spherical harmonic coefficients ($9 \times 3$ color channels).

Color computation:
\begin{equation}
    c(\mathbf{d}) = \text{sigmoid}\left(\sum_{l=0}^{2} \sum_{m=-l}^{l} k_l^m Y_l^m(\mathbf{d})\right)
\end{equation}

Volume rendering via trilinear interpolation of 8 nearest corners.

\subsubsection{Multi-Resolution and Sparsity}

\begin{itemize}
    \item Coarse-to-fine: $128^3 \to 256^3 \to 512^3$
    \item Periodic pruning: remove voxels with $\sigma < \epsilon$
    \item Typical occupancy: 5--20\% of total grid
    \item Native sparse storage
\end{itemize}

\subsubsection{RDD Compatibility}

\textbf{Excellent}:
\begin{itemize}
    \item Natural $N^3$ voxel grid with 28 channels per voxel
    \item Native sparsity (periodic pruning)
    \item Coarse-to-fine upsampling is standard practice
    \item \textbf{Challenge}: High memory at high resolution ($512^3 \times 28$ = 14GB dense)
    \item RDD's sparse approach directly solves this memory problem
\end{itemize}

%-------------------------------------------------------------------
\subsection{K-Planes Factorized Representation}

\textbf{Paper}: ``K-Planes: Explicit Radiance Fields in Space, Time, and Appearance'' (CVPR 2023)\\
\textbf{ArXiv}: \url{https://arxiv.org/abs/2301.10241}

\subsubsection{Mathematical Formulation}

3D represented as 3 planes with Hadamard (element-wise) product:
\begin{equation}
    f(\mathbf{q}) = f_{XY}(x,y) \odot f_{XZ}(x,z) \odot f_{YZ}(y,z)
\end{equation}

Multi-scale version with $S$ spatial scales:
\begin{equation}
    f(\mathbf{q}) = \text{concat}\left[\bigodot_{c \in \{XY,XZ,YZ\}} f_c^{(s)}(\mathbf{q})\right]_{s=1}^{S}
\end{equation}

\subsubsection{Key Property: Implicit Sparsity via Hadamard Product}

If \emph{any} plane has a zero at a given projection:
\begin{equation}
    f_{XY}(x,y) = 0 \quad \Rightarrow \quad f(\mathbf{q}) = 0 \quad \forall z
\end{equation}

This provides a natural sparsity mechanism without explicit sparse data structures!

\subsubsection{Conversion to $N^3$ Grid}

\begin{equation}
    \text{Grid}[i,j,k] = P_{XY}[i,j] \odot P_{XZ}[i,k] \odot P_{YZ}[j,k]
\end{equation}

\begin{itemize}
    \item No MLP needed --- purely explicit evaluation
    \item Memory: $3 \times S \times N^2 \times M$ (vs $N^3 \times C$ for dense voxel)
    \item Typical: $N=256, M=32, S=4$: $\approx$25M params $\approx$ 100MB
\end{itemize}

\subsubsection{RDD Compatibility}

\textbf{Very high}:
\begin{itemize}
    \item Upsampling planes (2D bilinear) directly equals volume upsampling
    \item Empty voxels detectable by checking plane values
    \item Constant upsampling on planes: $O(N^2)$ instead of $O(N^3)$
    \item Diffusion can operate on \textbf{plane factors} rather than full volume
\end{itemize}

\textbf{Novel idea}: Run RDD diffusion on the \emph{plane factors} directly. Multi-resolution: $64^2 \to 128^2 \to 256^2$ planes. The resulting 3D is implicitly multi-resolution.

%-------------------------------------------------------------------
\subsection{TensoRF VM Decomposition}

\textbf{Paper}: ``TensoRF: Tensorial Radiance Fields'' (ECCV 2022)\\
\textbf{ArXiv}: \url{https://arxiv.org/abs/2203.09517}

\subsubsection{Vector-Matrix Decomposition}

\begin{equation}
    \mathcal{T} = \sum_{r=1}^{R_1} \mathbf{v}_r^X \circ \mathbf{M}_r^{Y,Z} + \sum_{r=1}^{R_2} \mathbf{v}_r^Y \circ \mathbf{M}_r^{X,Z} + \sum_{r=1}^{R_3} \mathbf{v}_r^Z \circ \mathbf{M}_r^{X,Y}
\end{equation}

where $\mathbf{v}_r^X \in \mathbb{R}^I$ (vector) and $\mathbf{M}_r^{Y,Z} \in \mathbb{R}^{J \times K}$ (matrix/plane).

\begin{itemize}
    \item Memory: $O(N^2)$ instead of $O(N^3)$
    \item Quality: comparable to dense voxel at $300^3$
    \item Size: $<75$MB vs Plenoxels' GBs
\end{itemize}

\subsubsection{RDD Compatibility}

\textbf{High}:
\begin{itemize}
    \item Can evaluate to any $N^3$ grid for diffusion
    \item Factored structure allows efficient upsampling (2D plane upsample + 1D vector upsample)
    \item Compact enough for latent diffusion
\end{itemize}

%=================================================================
\section{Category 3: Grid-Based Implicit/Explicit 3D Generation}
%=================================================================

\subsection{SDFusion: Latent Diffusion on SDF Grids (CVPR 2023)}

\textbf{ArXiv}: \url{https://arxiv.org/abs/2212.04493}

\subsubsection{Method}

\begin{itemize}
    \item Input: Truncated SDF on $64^3$ voxel grid
    \item VQ-VAE compression: $64^3 \to 8^3$ latent ($8\times$ per axis)
    \item Latent diffusion: 3D U-Net denoiser on $8^3$ latent space
    \item Multi-modal conditioning: image (CLIP), text (CLIP), partial shape
\end{itemize}

\subsubsection{Limitations for RDD}

\begin{itemize}
    \item[$\times$] Dense 3D convolutions on full grid
    \item[$\times$] Fixed resolution $64^3$ (hard to scale)
    \item[$\times$] No sparsity exploitation
    \item[$\checkmark$] But: VQ-VAE + latent diffusion paradigm is sound
\end{itemize}

%-------------------------------------------------------------------
\subsection{DiffRF: Direct Voxel Radiance Field Diffusion (CVPR 2023)}

\textbf{ArXiv}: \url{https://arxiv.org/abs/2212.01206}

\subsubsection{Key Innovation: Rendering-Guided Loss}

\begin{equation}
    \mathcal{L} = \mathcal{L}_{\text{DDPM}} + \lambda \mathcal{L}_{\text{RGB}}
\end{equation}
\begin{equation}
    \mathcal{L}_{\text{RGB}} = \|\hat{I}(\hat{V}_0) - I_{\text{gt}}\|^2
\end{equation}

where $\hat{V}_0$ is the predicted clean volume from noisy input, and $\hat{I}$ is differentiable volume rendering.

\subsubsection{Architecture}

\begin{itemize}
    \item Direct diffusion on $128^3 \times 4$ voxel grid (density + RGB)
    \item 3D U-Net denoiser (converted from 2D ADM U-Net)
    \item No encoder/decoder --- operates in voxel space directly
\end{itemize}

\subsubsection{RDD Relevance}

Rendering-guided loss is valuable for RDD-3D:
\begin{itemize}
    \item Can supervise multi-resolution generation with 2D rendered images
    \item Avoids need for perfect 3D ground truth
    \item Compatible with any voxel-based representation
\end{itemize}

%-------------------------------------------------------------------
\subsection{VoxGRAF: Sparse Voxel GAN (NeurIPS 2022)}

\textbf{ArXiv}: \url{https://arxiv.org/abs/2206.07695}

\subsubsection{Architecture}

\begin{lstlisting}[style=flowchart, caption={VoxGRAF Progressive Growing}]
z (latent) --> Dense 32^3 feature grid
    |
    v (3D sparse conv)
Grow to 64^3 --> Free Space Pruning (discard empty)
    |
    v (3D sparse conv)  
Grow to 128^3 --> Free Space Pruning
    |
    v
Volume Rendering --> 2D image --> 2D Discriminator
\end{lstlisting}

\subsubsection{RDD Relevance}

VoxGRAF's \textbf{progressive growing + free-space pruning} is exactly the paradigm RDD enables for diffusion:
\begin{itemize}
    \item Start dense at low resolution
    \item Upsample and prune empty voxels
    \item Process only non-empty voxels at high resolution
    \item Replace GAN with RDD diffusion
\end{itemize}

%=================================================================
\section{Category 4: Mesh/Surface Extraction from Grids}
%=================================================================

\subsection{FlexiCubes (SIGGRAPH 2023)}

\textbf{ArXiv}: \url{https://arxiv.org/abs/2308.05371}

\subsubsection{Differentiable Mesh from Voxel Grid}

Given SDF values $s$ on a cubic grid, FlexiCubes extracts high-quality mesh with extra learnable parameters:

\textbf{Weight parameters} $\alpha \in \mathbb{R}_{>0}^8$ per cube: adjust edge crossing positions:
\begin{equation}
    u_e = \frac{s(x_i)\alpha_i \cdot x_j - s(x_j)\alpha_j \cdot x_i}{s(x_i)\alpha_i - s(x_j)\alpha_j}
\end{equation}

\textbf{Split parameters} $\beta \in [0,1]^{12}$ per cube: control tetrahedralization.

\subsubsection{RDD Integration}

FlexiCubes provides differentiable mesh extraction from any grid-based SDF:
\begin{itemize}
    \item RDD generates sparse SDF grid (multi-resolution)
    \item FlexiCubes converts final grid to mesh
    \item End-to-end differentiable for rendering losses
    \item Supports $128^3$ to $256^3$ grids
\end{itemize}

%-------------------------------------------------------------------
\subsection{DMTet: Deep Marching Tetrahedra (NeurIPS 2021)}

\textbf{ArXiv}: \url{https://arxiv.org/abs/2111.04276}

\subsubsection{Key Properties}

\begin{itemize}
    \item Tetrahedral grid (from regular cube subdivision)
    \item Each vertex: SDF value $s_i$ + deformation $\Delta v_i$
    \item Differentiable mesh extraction via linear interpolation on sign-changing edges
    \item \textbf{Volume subdivision}: only refine tetrahedra containing surface
    \item Resolution: $128^3$ base, up to $256^3$ with adaptive subdivision
\end{itemize}

\subsubsection{Surface point extraction}

\begin{equation}
    p_{ij} = \frac{|s_j| \cdot (v_i + \Delta v_i) + |s_i| \cdot (v_j + \Delta v_j)}{|s_i| + |s_j|}
\end{equation}

\subsubsection{RDD Integration}

DMTet's coarse-to-fine subdivision aligns with RDD:
\begin{itemize}
    \item Coarse grid ($32^3$): RDD generates rough SDF
    \item Adaptive subdivision: only refine near surface (= non-empty voxels)
    \item Fine grid ($128^3+$): RDD refines surface detail
\end{itemize}

%-------------------------------------------------------------------
\subsection{Get3D: SDF Generation on Grid (NeurIPS 2022)}

\textbf{ArXiv}: \url{https://arxiv.org/abs/2209.11163}

\subsubsection{Architecture}

\begin{itemize}
    \item StyleGAN2 backbone generates feature volume $4^3 \to 8^3 \to 16^3 \to 32^3 \times 64$
    \item Trilinear interpolation to DMTet vertices
    \item MLP predicts SDF + deformation at each vertex
    \item DMTet extracts mesh differentiably
    \item 2D discriminator on rendered images
\end{itemize}

\subsubsection{Grid Structure}

Feature volume progression:
\begin{equation}
    F^{(l)} = \text{ModConv3D}(F^{(l-1)}; w) \quad l = 1,...,4
\end{equation}
\begin{equation}
    F^{(0)} \in \mathbb{R}^{4\times4\times4\times256}, \quad F^{(4)} \in \mathbb{R}^{32\times32\times32\times64}
\end{equation}

\subsubsection{RDD Compatibility}

The progressive $4^3 \to 32^3$ feature volume is directly analogous to RDD's multi-resolution stages. Replace StyleGAN2 with RDD diffusion on the feature volume.

%=================================================================
\section{Category 5: NeRF-to-Grid Representations for Diffusion}
%=================================================================

\subsection{Tri-plane + Diffusion (EG3D / RodinDiffusion / NFD)}

\subsubsection{EG3D Tri-plane (CVPR 2022)}

Three orthogonal feature planes:
\begin{equation}
    F(\mathbf{p}) = T_{XY}(x,y) + T_{XZ}(x,z) + T_{YZ}(y,z), \quad T_* \in \mathbb{R}^{N \times N \times C}
\end{equation}

Decoded via MLP: $(\sigma, \mathbf{c}) = \text{MLP}(F(\mathbf{p}))$

\subsubsection{Diffusion on Tri-planes}

Multiple works apply diffusion directly on tri-plane latents:
\begin{itemize}
    \item \textbf{RodinDiffusion} (CVPR 2023): 3D avatar generation, $512^2 \times 64$ planes
    \item \textbf{NFD} (CVPR 2023): reshape tri-planes as 2D images, use 2D U-Net diffusion
    \item \textbf{SSDNeRF} (ICCV 2023): joint optimization of tri-planes + diffusion prior
\end{itemize}

\subsubsection{RDD Compatibility}

\textbf{Moderate}:
\begin{itemize}
    \item Tri-planes are 2D, not 3D --- need architectural adaptation
    \item Cannot directly identify ``empty voxels'' without MLP evaluation
    \item Memory efficient ($O(N^2)$) but lacks native 3D sparsity
    \item \textbf{Alternative}: Run RDD on tri-plane factors (2D multi-resolution diffusion $\times$ 3 planes)
\end{itemize}

%-------------------------------------------------------------------
\subsection{DVGO / Direct Voxel Grid Optimization (CVPR 2022)}

\subsubsection{Post-Activation Interpolation}

Key innovation for sharp boundaries:
\begin{equation}
    \sigma(\mathbf{x}) = \text{softplus}\left(\text{trilinear}(V^{(\text{density})}, \mathbf{x}) + b\right)
\end{equation}

Store raw (pre-activation) values; interpolate first, activate after.

\subsubsection{Grid Structure}

\begin{itemize}
    \item Density grid: $V^{(\text{density})} \in \mathbb{R}^{N \times N \times N}$
    \item Feature grid: $V^{(\text{feature})} \in \mathbb{R}^{N \times N \times N \times C}$ ($C=12$)
    \item Shallow MLP decodes features to color
    \item Resolution: $160^3$ to $320^3$
\end{itemize}

\subsubsection{RDD Compatibility}

\textbf{Good}: DVGO's density grid provides natural sparsity mask. Feature grid can be ``baked'' (evaluate MLP at grid points) to produce pure grid representation suitable for diffusion.

%=================================================================
\section{Detailed Comparison Table}
%=================================================================

\begin{longtable}{p{2.5cm}|p{2cm}|p{1.5cm}|p{1.5cm}|p{1.5cm}|p{2cm}|p{2cm}}
\caption{Comprehensive Comparison of 3D Representations for Sparse RDD} \\
\toprule
\textbf{Method} & \textbf{Representation} & \textbf{$N^3$ Grid} & \textbf{Sparse} & \textbf{Multi-Res} & \textbf{RDD Compat.} & \textbf{Max Res.} \\
\midrule
\endfirsthead
\toprule
\textbf{Method} & \textbf{Representation} & \textbf{$N^3$ Grid} & \textbf{Sparse} & \textbf{Multi-Res} & \textbf{RDD Compat.} & \textbf{Max Res.} \\
\midrule
\endhead
XCube & Sparse VDB & \checkmark & \checkmark\checkmark & \checkmark\checkmark & $\star\star\star\star\star$ & 1024$^3$ \\
\midrule
OctFusion & Octree latent & \checkmark & \checkmark\checkmark & \checkmark\checkmark & $\star\star\star\star\star$ & 256$^3$ \\
\midrule
TRELLIS & Sparse SLAT & \checkmark & \checkmark\checkmark & \checkmark & $\star\star\star\star\star$ & 1536$^3$ \\
\midrule
\textbf{GaussianCube} & \textbf{OT-structured 3DGS} & \checkmark\checkmark & \checkmark & \checkmark & $\star\star\star\star\star$ & \textbf{32$^3$ (extensible)} \\
\midrule
Direct3D-S2 & Sparse DiT & \checkmark & \checkmark\checkmark & --- & $\star\star\star\star$ & 1024$^3$ \\
\midrule
VoxGRAF & Sparse voxel & \checkmark & \checkmark\checkmark & \checkmark & $\star\star\star\star$ & 128$^3$ \\
\midrule
Plenoxels & Dense/sparse SH & \checkmark\checkmark & \checkmark & \checkmark & $\star\star\star\star$ & 512$^3$ \\
\midrule
K-Planes & Factored planes & indirect & implicit & \checkmark & $\star\star\star\star$ & 512$^3$ \\
\midrule
TensoRF-VM & Vec-Mat decomp & indirect & implicit & \checkmark & $\star\star\star\star$ & 300$^3$ \\
\midrule
Get3D & Dense feature vol & \checkmark & --- & \checkmark & $\star\star\star$ & 32$^3$ feat \\
\midrule
SDFusion & Dense SDF grid & \checkmark & --- & --- & $\star\star$ & 64$^3$ \\
\midrule
DiffRF & Dense voxel RF & \checkmark & --- & --- & $\star\star$ & 128$^3$ \\
\midrule
3DShape2VecSet & Token set & --- & --- & --- & $\star\star$ & N/A \\
\midrule
Tri-plane & 3$\times N^2$ planes & indirect & --- & --- & $\star\star$ & 256$^2$ \\
\midrule
MeshGPT & 1D sequence & --- & --- & --- & $\star$ & N/A \\
\midrule
Shap-E & MLP params & --- & --- & --- & $\star$ & N/A \\
\bottomrule
\end{longtable}

%=================================================================
\section{Proposed Architecture: RDD-3D Sparse Diffusion}
%=================================================================

\subsection{Design Choices}

Based on this survey, the optimal architecture combines:

\begin{enumerate}
    \item \textbf{Representation}: GaussianCube (proven quality, native grid) OR Sparse Voxel SDF (XCube-style, native sparsity)
    \item \textbf{Encoder}: Sparse 3D VAE (TRELLIS-style shifted-window attention)
    \item \textbf{Diffusion}: RDD with unified multi-scale denoiser (OctFusion philosophy)
    \item \textbf{Backbone}: Sparse 3D U-Net (XCube) or Sparse DiT with SSA (Direct3D-S2)
    \item \textbf{Decoder}: Differentiable rendering (GaussianCube splatting or FlexiCubes mesh extraction)
\end{enumerate}

\subsection{Architecture Option A: RDD-GaussianCube}

\begin{lstlisting}[style=flowchart, caption={RDD-GaussianCube Architecture}]
=== TRAINING ===

    3D Asset (mesh/scan)
        |
    3DGS Fitting (density-constrained, N_max=64^3)
        |
    Optimal Transport Structuralization
        |
    GaussianCube: 64 x 64 x 64 x 14
        |
    Sparse VAE Encoder (3D sparse conv)
        |-- Downsample: 64^3 -> 32^3 -> 16^3 -> 8^3
        |-- KL regularization
        |
    Latent: 8^3 x D (dense at coarsest level)
        |
    RDD Forward Process:
        |-- Stage 1 (t in [T, T2]): drift toward 4^3 average
        |-- Stage 2 (t in [T2, T1]): drift toward 8^3 
        |-- Stage 3 (t in [T1, 0]): no drift (finest detail)
        |
    Train Sparse 3D U-Net denoiser

=== GENERATION ===

    Sample z_T ~ N(0, I) at 8^3 x D
        |
    RDD Reverse (Stage 1): 8^3 -> clean coarse structure
        |
    Constant Upsample: 8^3 -> 16^3 (nearest-neighbor)
        |-- Threshold: keep only voxels with predicted density > eps
        |-- Result: sparse 16^3 
        |
    RDD Reverse (Stage 2): denoise sparse 16^3
        |
    Constant Upsample: 16^3 -> 32^3 -> 64^3 (with pruning)
        |
    RDD Reverse (Stage 3): denoise sparse 64^3 (fine detail)
        |
    Sparse VAE Decoder -> GaussianCube 64^3 x 14
        |
    Gaussian Splatting Rendering -> Novel views
\end{lstlisting}

\subsection{Architecture Option B: RDD-SparseVoxel (XCube-inspired)}

\begin{lstlisting}[style=flowchart, caption={RDD-SparseVoxel Architecture}]
=== TRAINING ===

    3D Asset -> Voxelize to Sparse SDF (OpenVDB)
        |
    Multi-level sparse representation:
        Level 1: 16^3 (dense, coarse structure)
        Level 2: 64^3 (sparse, medium detail)
        Level 3: 256^3 (sparse, fine surface)
        |
    Sparse Structure VAE (per level):
        Encode (sparse 3D conv) -> latent per level
        |
    RDD Forward (UNIFIED, single chain):
        x_t = sqrt(alpha_bar_t) * x_0 + sqrt(alpha_bar_t)*kappa(t)*d + noise
        |
        d = UP(DOWN(x_0)) - x_0   [constant upsample!]
        |
        Single noise schedule covers ALL levels
        kappa(t) smoothly transitions between resolution stages
        |
    Train UNIFIED Sparse 3D U-Net (OctFusion-style weight sharing)

=== GENERATION ===

    z_T ~ N(0, I) at Level 1 (16^3 dense)
        |
    RDD Reverse (t: T -> T2):
        Denoises toward coarse structure
        kappa: 1 -> 0 (drift disappears)
        |
    [TRANSITION: constant upsample 16^3 -> 64^3]
        |-- Mark empty voxels (density < eps)
        |-- Discard empty -> sparse 64^3
        |
    RDD Reverse (t: T2 -> T1):
        Denoises sparse 64^3
        |
    [TRANSITION: constant upsample 64^3 -> 256^3]
        |-- Further prune empty
        |
    RDD Reverse (t: T1 -> 0):
        Final surface detail at sparse 256^3
        |
    Decode: latent -> SDF -> FlexiCubes -> Mesh
\end{lstlisting}

\subsection{Architecture Option C: RDD-FactoredPlanes (Novel)}

\begin{lstlisting}[style=flowchart, caption={RDD on Factored Plane Representation}]
=== KEY IDEA ===
    Instead of diffusing on N^3 volume, diffuse on 
    3 planes (each N^2 x C) using K-Planes factorization.
    
    Memory: O(N^2) instead of O(N^3)!
    RDD upsampling: 2D plane upsample (constant/nearest)
    Sparsity: Hadamard product guarantees zero propagation

=== PIPELINE ===

    3D Asset -> Fit K-Planes (multi-scale)
        P_XY^s, P_XZ^s, P_YZ^s for s=1,...,S
        |
    Stack/concatenate planes -> "triplane latent"
        Total: 3 x S x N^2 x M
        |
    RDD 2D Diffusion on each plane independently:
        Coarse planes (64^2): diffuse first
        Fine planes (256^2): diffuse conditioned on coarse
        |
    Constant upsample on 2D planes (nearest-neighbor)
        |
    Reconstruct: f(x,y,z) = Product of plane features
        |
    Threshold density -> Marching Cubes -> Mesh

=== ADVANTAGES ===
    - O(N^2) memory and compute (not N^3!)
    - Can reuse existing 2D diffusion architectures
    - RDD constant upsample works in 2D
    - Implicit 3D sparsity from Hadamard product
    
=== DISADVANTAGES ===
    - Cross-plane artifacts (known tri-plane issue)
    - Cannot directly identify sparse 3D regions from 2D planes
    - Less natural for 3D sparse operations
\end{lstlisting}

%=================================================================
\section{Experimental Baselines and Expected Metrics}
%=================================================================

\subsection{Datasets}

\begin{itemize}
    \item \textbf{ShapeNet} (55 categories, ~51K shapes): Standard benchmark for 3D generation
    \item \textbf{Objaverse} (800K+ objects): Large-scale diverse 3D dataset
    \item \textbf{OmniObject3D} (6K real scanned objects, 190 categories): Real-world benchmark
    \item \textbf{Waymo Open Dataset}: Large-scale outdoor scenes (for scene-level generation)
\end{itemize}

\subsection{Metrics}

\begin{itemize}
    \item \textbf{FID} (Frechet Inception Distance): on multi-view rendered images (20 views standard)
    \item \textbf{KID} (Kernel Inception Distance): unbiased alternative to FID
    \item \textbf{1-NNA} (1-Nearest Neighbor Accuracy): on Chamfer Distance (CD) and Earth Mover's Distance (EMD)
    \item \textbf{Coverage (COV)}: diversity measure
    \item \textbf{MMD} (Minimum Matching Distance): fidelity measure
    \item \textbf{PSNR/SSIM/LPIPS}: for novel-view synthesis quality
    \item \textbf{IoU}: volumetric overlap for occupancy evaluation
    \item \textbf{Chamfer Distance}: point cloud geometry accuracy
\end{itemize}

\subsection{Expected Performance Targets}

Based on current SOTA (XCube, TRELLIS, GaussianCube):

\begin{table}[H]
\centering
\caption{Target Metrics for RDD-3D}
\begin{tabular}{lcc}
\toprule
Metric & Current SOTA & RDD-3D Target \\
\midrule
ShapeNet FID (multi-view) & $\sim$15--25 & $\leq$ 20 \\
Generation time (256$^3$) & 2.5s (OctFusion) & $\leq$ 3s \\
Memory (256$^3$) & 8GB (dense) & $\leq$ 2GB (sparse) \\
Max resolution & 1536$^3$ (TRELLIS) & 512$^3$+ \\
\bottomrule
\end{tabular}
\end{table}

%=================================================================
\section{Action Plan}
%=================================================================

\subsection{P0: Proof of Concept (4--6 weeks)}

\begin{enumerate}
    \item \textbf{Implement RDD-GaussianCube baseline}:
    \begin{itemize}
        \item Use GaussianCube's public code for data preparation
        \item Implement 3D RDD forward process (constant upsample + drift)
        \item Train 3D U-Net denoiser at $32^3$ resolution
        \item Validate: compare FID with GaussianCube's original ADM baseline
    \end{itemize}
    
    \item \textbf{Implement sparse masking}:
    \begin{itemize}
        \item After coarse-level generation, threshold opacity to create binary mask
        \item Verify: generation quality preserved when discarding low-opacity voxels
        \item Measure: actual sparsity ratio achieved (target: 70--90\% empty)
    \end{itemize}
    
    \item \textbf{Benchmark constant vs. bilinear upsampling}:
    \begin{itemize}
        \item Ablation: nearest-neighbor vs bilinear upsample in RDD 3D
        \item Confirm: constant upsample + RDD drift correction matches bilinear quality
    \end{itemize}
\end{enumerate}

\subsection{P1: Sparse Acceleration (4--6 weeks)}

\begin{enumerate}
    \item \textbf{Implement sparse 3D U-Net}:
    \begin{itemize}
        \item Replace dense 3D convolutions with sparse alternatives (MinkowskiEngine or fVDB)
        \item Only process non-empty voxels at medium/fine resolutions
        \item Measure: FLOPs reduction and wallclock speedup
    \end{itemize}
    
    \item \textbf{Scale to higher resolution}:
    \begin{itemize}
        \item Extend from $32^3$ to $64^3$ and $128^3$ with sparse operations
        \item Implement multi-stage RDD: $16^3 \to 32^3 \to 64^3 \to 128^3$
        \item Verify: linear scaling with occupied voxels (not cubic with resolution)
    \end{itemize}
    
    \item \textbf{Sparse attention (optional)}:
    \begin{itemize}
        \item Implement Direct3D-S2's Spatial Sparse Attention
        \item Compare: sparse U-Net vs sparse DiT backbone
    \end{itemize}
\end{enumerate}

\subsection{P2: Full System and Paper (6--8 weeks)}

\begin{enumerate}
    \item \textbf{Alternative representations}:
    \begin{itemize}
        \item Test K-Planes factored representation with RDD
        \item Compare: GaussianCube vs. SDF voxel vs. factored planes
        \item Identify best representation for different tasks (objects vs. scenes)
    \end{itemize}
    
    \item \textbf{Conditioning and applications}:
    \begin{itemize}
        \item Add image conditioning (CLIP/DINOv2)
        \item Add text conditioning
        \item Shape completion from partial input
    \end{itemize}
    
    \item \textbf{Paper writing}:
    \begin{itemize}
        \item Core contribution: ``Sparse Multi-Resolution 3D Diffusion via Constant Upsampling''
        \item Key claim: RDD's neighbor-independent upsampling enables native sparse 3D generation
        \item Ablation: $\gamma=0$ (no drift) vs full RDD vs cascade baseline
        \item Comparison: vs XCube (cascade), vs OctFusion (unified but no sparsity claim), vs TRELLIS (two-stage)
    \end{itemize}
\end{enumerate}

%=================================================================
\section{Key Novelty Claim for Paper}
%=================================================================

\subsection{Positioning Statement}

\begin{center}
\fbox{\parbox{0.92\textwidth}{
\textbf{Novelty}: We observe that RDD's constant (nearest-neighbor) upsampling creates a unique opportunity for 3D generation: since upsampled values are \emph{independent of neighbors}, empty voxels can be identified and \emph{discarded without loss of information}. This transforms the $O(N^3)$ dense volumetric diffusion into $O(k)$ sparse generation where $k$ is the number of surface voxels ($k \propto N^2 \ll N^3$).

\textbf{Unlike prior work}:
\begin{itemize}
    \item XCube uses \emph{cascade} diffusion (separate models per level) --- we use a \emph{single unified RDD chain}
    \item OctFusion handles multi-scale but doesn't exploit the \emph{mathematical guarantee} that constant upsampling enables exact sparse computation
    \item TRELLIS separates structure and content generation --- RDD's drift term naturally handles both in one process
\end{itemize}

\textbf{Mathematical insight}: In RDD, if a block of voxels has identical values (constant region), the drift term $\mathbf{d} = \text{UP}(\text{DOWN}(x_0)) - x_0 = 0$ exactly. Empty (zero) regions remain zero throughout the forward/reverse process, making sparse computation \emph{exact} (not approximate).
}}
\end{center}

\subsection{Comparison with Existing Work}

\begin{table}[H]
\centering
\caption{Novelty Positioning}
\begin{tabular}{p{3cm}cccc}
\toprule
& Single Chain & Constant UP & Exact Sparsity & Unified Model \\
\midrule
XCube & $\times$ (cascade) & N/A & $\times$ & $\times$ \\
OctFusion & \checkmark & $\times$ (learned) & $\times$ (approx) & \checkmark \\
TRELLIS & $\times$ (two-stage) & N/A & $\times$ & $\times$ \\
Cascade DDPM & $\times$ & $\times$ (bilinear) & $\times$ & $\times$ \\
\textbf{RDD-3D (ours)} & \checkmark & \checkmark & \checkmark & \checkmark \\
\bottomrule
\end{tabular}
\end{table}

%=================================================================
\section{References}
%=================================================================

\begin{enumerate}[label={[\arabic*]}]
    \item XCube: Ren et al., ``XCube: Large-Scale 3D Generative Modeling using Sparse Voxel Hierarchies,'' CVPR 2024 Highlight. arXiv:2312.03806
    \item OctFusion: Xiong et al., ``OctFusion: Octree-based Diffusion Models for 3D Shape Generation,'' 2024. arXiv:2408.14732
    \item TRELLIS: Xiang et al., ``Structured 3D Latents for Scalable and Versatile 3D Generation,'' 2024. arXiv:2412.01506
    \item GaussianCube: Zhang et al., ``GaussianCube: Structuring Gaussian Splatting using Optimal Transport for 3D Generative Modeling,'' NeurIPS 2024. arXiv:2403.19655
    \item Direct3D: Wu et al., ``Direct3D: Scalable Image-to-3D Generation via 3D Latent Diffusion Transformer,'' 2024. arXiv:2405.14832
    \item Direct3D-S2: ``Gigascale 3D Generation Made Easy with Spatial Sparse Attention,'' 2025. arXiv:2505.17412
    \item 3DShape2VecSet: Zhang et al., ``3DShape2VecSet: A 3D Shape Representation for Neural Fields and Generative Diffusion Models,'' SIGGRAPH 2023. arXiv:2301.11445
    \item Michelangelo: Zhao et al., ``Michelangelo: Conditional 3D Shape Generation with Shape-Image-Text Aligned Latent Representation,'' NeurIPS 2023. arXiv:2306.17115
    \item SDFusion: Cheng et al., ``SDFusion: Multimodal 3D Shape Completion, Reconstruction, and Generation,'' CVPR 2023. arXiv:2212.04493
    \item DiffRF: M\"uller et al., ``DiffRF: Rendering-Guided 3D Radiance Field Diffusion,'' CVPR 2023. arXiv:2212.01206
    \item VoxGRAF: Schwarz et al., ``VoxGRAF: Fast 3D-Aware Image Synthesis with Sparse Voxel Grids,'' NeurIPS 2022. arXiv:2206.07695
    \item SSDNeRF: Chen et al., ``Single-Stage Diffusion NeRF: A Unified Approach to 3D Generation and Reconstruction,'' ICCV 2023. arXiv:2304.06714
    \item TensoRF: Chen et al., ``TensoRF: Tensorial Radiance Fields,'' ECCV 2022. arXiv:2203.09517
    \item K-Planes: Fridovich-Keil et al., ``K-Planes: Explicit Radiance Fields in Space, Time, and Appearance,'' CVPR 2023. arXiv:2301.10241
    \item Plenoxels: Fridovich-Keil et al., ``Plenoxels: Radiance Fields without Neural Networks,'' CVPR 2022. arXiv:2112.05131
    \item DVGO: Sun et al., ``Direct Voxel Grid Optimization: Super-fast Convergence for Radiance Fields Reconstruction,'' CVPR 2022. arXiv:2111.11215
    \item EG3D: Chan et al., ``Efficient Geometry-aware 3D Generative Adversarial Networks,'' CVPR 2022. arXiv:2112.07945
    \item FlexiCubes: Shen et al., ``Flexible Isosurface Extraction for Gradient-Based Mesh Optimization,'' SIGGRAPH 2023. arXiv:2308.05371
    \item DMTet: Shen et al., ``Deep Marching Tetrahedra: a Hybrid Representation for High-Resolution 3D Shape Synthesis,'' NeurIPS 2021. arXiv:2111.04276
    \item Get3D: Gao et al., ``GET3D: A Generative Model of High Quality 3D Textured Shapes Learned from Images,'' NeurIPS 2022. arXiv:2209.11163
    \item NKSR: Huang et al., ``Neural Kernel Surface Reconstruction,'' CVPR 2023. arXiv:2305.19590
    \item LION: Zeng et al., ``LION: Latent Point Diffusion Models for 3D Shape Generation,'' NeurIPS 2022.
    \item MeshGPT: Siddiqui et al., ``MeshGPT: Generating Triangle Meshes with Decoder-Only Transformers,'' ICLR 2024. arXiv:2311.15475
    \item Shap-E: Jun and Nichol, ``Shap-E: Generating Conditional 3D Implicit Functions,'' 2023. arXiv:2305.02463
    \item PVD: Zhou et al., ``3D Shape Generation and Completion through Point-Voxel Diffusion,'' ICCV 2021. arXiv:2104.03670
    \item Conv. Occ. Networks: Peng et al., ``Convolutional Occupancy Networks,'' ECCV 2020.
    \item L3DG: Lan et al., ``Latent 3D Gaussian Diffusion,'' 2024. arXiv:2410.13530
    \item SPVD: ``Sparse Point-Voxel Diffusion,'' 2024. arXiv:2408.06145
    \item fVDB: Williams et al., ``fVDB: A Deep-Learning Framework for Sparse, Large-Scale, and High-Performance Spatial Intelligence,'' 2024. arXiv:2407.01781
\end{enumerate}

\end{document}
